\chapter{Adatbevitel: typed SQL, CRUD, formok és PRG}

Egy klasszikus webalkalmazás alapciklusa: adatot olvasunk az adatbázisból, HTML formon typed inputot fogadunk, validálunk, autorizáljuk az objectet vagy műveletet, tranzakcióban módosítunk, majd redirect után újra GET-tel olvasunk. Ebben a fejezetben ezt a teljes ciklust egy helyen kezeljük. A fontos különbség: a \emph{valid input még nem authorized input}.

\section{Typed SQL és CRUD}

\subsection*{Lekérdezés}

\begin{lstlisting}[language=RWLang]
query fn articleBySlug(db: Db, slug: Slug)
    -> Result<Article, DbError> sql {
    SELECT id, slug, title, body
    FROM articles
    WHERE slug = :slug
}
\end{lstlisting}

Az SQL szöveg statikus és compiler-owned. Felhasználói érték csak typed bindként, például \code{:slug}, kerülhet bele.

\subsection*{Mutáció transactionben}

\begin{lstlisting}[language=RWLang]
query fn articleById(db: Db, id: Int)
    -> Result<Article, DbError> sql {
    SELECT id, slug, title, body
    FROM articles
    WHERE id = :id
}

query fn updateArticle(
    tx: Transaction,
    id: Int,
    title: String,
    body: String
) -> Result<Void, DbError> mutates Article by id sql {
    UPDATE articles
    SET title = :title, body = :body
    WHERE id = :id
}
\end{lstlisting}

Az action betölti az objectet, létrehozza az authorization proofot, és csak ezután használja az autorizált object kulcsát a mutációhoz:

\begin{lstlisting}[language=RWLang]
action fn update(
    ctx: ActionContext,
    db: Db,
    id: Int,
    title: String,
    body: String
) -> Result<Redirect, PageError> {
    let article = articleById(db, id)?;
    authorize article authenticated;
    transaction db {
        updateArticle(tx, article.id, title, body)?;
    }
    return Ok(redirect(adminArticles()));
}
\end{lstlisting}

Hétköznapi lokális tranzakciónál siker esetén commit, ismert adatbázishibánál rollback történik. A \code{Transaction} capability nem escape-elhet a blokkból. Retry-sensitive critical műveletnél viszont explicit transaction outcome-ot kell kezelni ahelyett, hogy minden hibáról rollbacket feltételeznénk; a \code{CommitUnknown} szemantikát a \ref{ch:db-deep-dive-hu}. fejezet részletezi.

\subsection*{Portable adatbázis-séma}

A V1 SQLite, PostgreSQL és MariaDB irányban is hordozható mintát céloz. \code{Date}, \code{DateTime}, \code{Uuid} és \code{Decimal} portable reprezentációja canonical text lehet.

\subsection*{Migration}

A szerver nem migrál automatikusan startupkor. Ajánlott folyamat:

\begin{lstlisting}
migrate verify
migrate apply
migrate verify
\end{lstlisting}

A migration credential legyen elkülönítve a normál application runtime credentialtől.


\section{Formok, validáció és PRG}

\subsection*{Named form schema}

\begin{lstlisting}[language=RWLang]
form ArticleForm {
    title<String>
    body<String>
    published<Bool>
    validate title length 2 200 body length 1 200000
}
\end{lstlisting}

Route:

\begin{lstlisting}[language=RWLang]
route articleCreate POST "/admin/articles"
    form ArticleForm
    auth role Editor
    => create;
\end{lstlisting}

A handler paramétersorrendje és típusa kövesse a form mezőit.

\subsection*{Hibás beküldés}

Named formnál a hiányzó kötelező mező, a mezőszintű typed decode hiba vagy a deklarált validation hibája 422 \emph{Unprocessable Content}. A runtime a mezőértékeket escape-elve tölti vissza, mezőhibát ad és CSRF tokent illeszt a formba. Ismeretlen vagy duplikált mező viszont 400, mert az request-schema megsértése.

Ez szándékosan más, mint az inline form- vagy JSON-schema: ott a typed decode/validation hibája request-contract sértésként \code{400}. Named formot akkor használj, amikor a hibát ugyanazon felhasználói űrlapon, mezőszintű visszajelzéssel akarod kezelni.

\subsection*{Sikeres POST és PRG}

Sikeres módosítás után redirectet adj vissza; így a böngésző GET-tel tölti be a céloldalt. Ez a Post/Redirect/Get minta megszünteti a véletlen újraküldés tipikus problémáját.

\subsection*{Egyszerű teljes form-action minta}

\begin{lstlisting}[language=RWLang]
action fn create(
    ctx: ActionContext,
    db: Db,
    title: String,
    body: String,
    published: Bool
) -> Result<Redirect, PageError> {
    let articleSlug = slug(title);
    transaction db {
        insertArticle(tx, articleSlug, title, body)?;
    }
    return Ok(redirect(adminArticles()));
}
\end{lstlisting}

A checkbox-szerű hiányzó \code{Bool} mező named form esetén \code{false}.

\subsection*{Gyakorlati döntési sorrend}

Állapotváltoztató form tervezésénél és review-jánál ezt a sorrendet használd:

\begin{enumerate}
  \item decode és validation a route/form boundaryn;
  \item meglévő rekord módosításánál töltsd be a cél objectet;
  \item hozd létre az object/permission/tenant authorization proofot;
  \item a megfelelő transaction vagy critical-operation contract alatt módosíts;
  \item typed GET route helperre redirectelj;
  \item csak azt a cache-t invalidáld, amelyet a módosítás ténylegesen elavulttá tett.
\end{enumerate}

Ez a sorrend kizár egy gyakori hibát: a jól formázott identifier vagy form önmagában nem bizonyítja, hogy a hívó módosíthatja a hivatkozott objectet.
